\chapter{Vector Fields}

\section{Vector Fields on Manifolds}
\begin{definition}[Section of a map]
   Let $\pi: M \to N$ be a continuous map, a \textbf{section of $\pi$} is a continuous right inverse for $\pi$.
\end{definition}
\begin{definition}
    Let $M$ be a smooth manifold, a \textbf{vector field on $M$} is a section $X: TM \to M$ of the map $\pi: TM \to M$. Generally, we call $X$ a \textbf{rough vector field on $M$} if its continuity is not given. The set of all smooth vector fields of $M$ is denoted by $\mathcal{X}(M)$
\end{definition}
Let $(U,x^i)$ be a local chart, as the case of tangent vectors, a rough vector field can be represented locally in terms of basis vector:
\[
    X_p = X^i(p) \frac{\pa}{\pa x^i}\bigg|_{p}
\]
Each function $X^i: U \to \bb{R}$ is called the \textbf{component function of $X$} in the given chart.
\begin{proposition}
    Let $X: M \to TM$ be a rough vector field. If $(U,x^i)$ is any smooth coordinate chart on $M$, the restriction $X: U \to TM$ is smooth if and only if its component functions reprect to this chart is smooth.
\end{proposition}
\begin{proof}
    Let $(x^i,v^i)$ be the natural coordinae on $\pi^{-1}(U) \subseteq TM$ associated with the chart $(U,x^i)$, then the coordinate representation of $X: M \to TM$ on $U$ is 
    \[
        \hat{X}(x) = (x^i,X^i(x))
    \]
    It follows immediately that the smoothness of $X$ on $U$ is equivalent to smoothness of its component functions.
\end{proof}
\begin{example}[Coordinate Vector Fields] Let $(U,x^i)$ be any smooth chart on $M$, the assignment 
    \[
        p \mapsto \frac{\pa}{\pa x^i}\bigg|_p
    \]
    determines a vector field on $U$.
\end{example}

\begin{example}[The Euler Vector Field] The vector field $V: \bb{R}^n \to T\bb{R}^n$ defined by 
    \[
        V(x) = x^1\frac{\pa}{\pa x^1}\bigg|_p + \dots + x^n\frac{\pa}{\pa x^n}\bigg|_p
    \]
is called the \textbf{the Euler vector field}.
    
\end{example}
\begin{lemma}[Extension Lemma for Vector Fields]
    Let $M$ be a smooth manifold, let $A \subseteq M$ be a closed subset. Suppose $X$ is a smooth vector field along $A$. Given any open subset $U$ containing $A$, there exists a smooth global vector field $\hat{X}$ on $M$ such that $\hat{X}|_A = X$ and $\supp \hat{X} \subseteq U$.
\end{lemma}
\begin{proof}
    Use extension lemma for smooth function.
\end{proof}
\begin{proposition}
    Let $M$ be a smooth manifold, then $\mc{X}(M)$ is a vector space.
\end{proposition}
\section{Local and Global Frames}
\begin{definition}
    Let $M^n$ be a smooth manifold and let $U \subseteq M$ be an open set. A collection of vector fields $\{E_i\}_{1 \leq i \leq n}$ called a \textbf{local frame} if for every point $p \in U$, the collection
    \[
        \{E_1|_p,\dots, E_n|_p\}
    \]
    forms a basis for $T_pM$. It is called a \textbf{global frame} if $U = M$.
\end{definition}
\begin{lemma}[Gram-Schmidt Algorithm for Frames] Let $(X_j)$ be a smooth local frame of $T\bb{R}^n$ over an open subset $U \subseteq \bb{R}^n$. Then there exists a smooth orthogonal frame $(E_j)$ over $U$ such that $\spann\{E_1|_p,\dots,E_j|_p\} = \spann \{X_1|_p,X_j|_p\}$ for each $j = 1,\dots, n$ and all $p \in U$.

\end{lemma}
\section{Vector Fields as Derivations}
\begin{proposition}
    Let $M$ be a smooth manifold and let $X: T \to TM$ be a rough vector field. The following are equivalent:
    \begin{enumerate}
        \item $X$ is smooth.
        \item $Xf$ is smooth on $M$ for all $f \in C^\infty(M)$.
        \item For every open subset $U \subseteq M$ and $f  \in C^\infty(M)$, the function $Xf$ is smooth on $U$.
    \end{enumerate}
\end{proposition}
\begin{proof}
    $(1 \ra 2)$. Suppose $X$ is smooth, let $(U,x^i)$ be any smooth coordinate on $U$ centered at $p$. Then for $x \in U$, one can write
    \[
        Xf(x) = \left(X^i(x)\frac{\pa}{\pa x^i}\bigg|_x\right)f = X^i(x)\frac{\pa f}{\pa x^i}(x)
    \]
    Since each component of $X$ is smooth, it follows that $Xf$ is smooth.

    $(2 \ra 3)$ follows since smoothness is a local property. Finally, to prove $(3 \ra 1)$, let $(U,x^i)$ be a smooth coordinate chart, for each $i = 1,\dots, n$ the function $X x^i X^i$ is smooth on $U$. Since each component of $X$ is smooth on $U$, then $X$ is smooth on $U$, implies that it is smooth on $M$.
    
\end{proof}
\begin{proposition}
    The vector field $X$ on $M$, in fact, is a derivation.
\end{proposition}
\begin{proposition}
    Let $M$ be a smooth manifold. A map $D: C^\infty(M) \to C^\infty(M)$ is a derivation if and only if it is of the form $Df = Xf$ for some smooth vector field $X \in \mc{X}(M)$.
\end{proposition}
\section{Pushforwarding Vector Fields}
\begin{definition}
    Let $F: M^m \to N^n$ be a smooth map between smooth manifolds. We say two vector field $X \in \mc{X}(M)$ and $Y \in \mc{X}(N)$ are \textbf{F-related} if $dF_p(X_p) = Y_{F(p)}$ for all $p \in M$.
\end{definition}
\begin{proposition}
    Let $F: M \to N$ be a smooth map and let $X \in \mc{X}(M), Y \in \mc{N}$. Then $X$ and $Y$ are $F$-related if and only if for all smooth real-valued function $f$ defined on an open subset of $N$, 
    \[
        X(f \circ F) = (Yf)\circ F.
    \]
\end{proposition}
\begin{proof}
    For any $p \in M$ and any smooth real-valued function $f$ defined on a neighborhood of $F(p)$,
    \[
        X(f \circ F)(p) = X_p (f \circ F) = dF_p(X_p)f,
    \]
    and 
    \[
        (Y f) \circ F(p)= Y_{F(p)}f
    \]
    Thus, this is true for all $f$ if and only if $X$ and $Y$ are $F$-related.
\end{proof}
\begin{proposition}
    Let $M$ and $N$ are smooth manifolds, and $F: M \to N$ is a diffeomorphism. For every $X \in \mc{X}(M)$, there is a unique smooth vector field on $N$ that is $F$-related to $X$.
\end{proposition}
\begin{proof}
    Since $F$ is a diffeomorphism, then one can define $Y$ by 
    \[
        Y_p = dF_{F^{-1}(p)}(X_{F^{-1}(p)}).
    \]
    Then $Y$ is a (rough) vector field, and since we have $Y = dF \circ X \circ F^{-1}$, it follows that $Y$ is smooth.
\end{proof}
\begin{definition}
    Let $F: M \to N$ be a diffeomorphism. The \textbf{pushforward of $X$ by $F$}, denoted by $F_*X$, is a unique vector field that is $F$-related to $X$, and can be defined explicitly by the formula
    \[
    \boxed{
    (F_*X)_p = dF_{F^{-1}(p)}(X_{F^{-1}(p)})}
    \]
\end{definition}
\begin{example}
    Let $M$ and $N$ be open submanifolds of $\bb{R}^2$ such that 
    \begin{align*}
        M &= \{(x,y) \mid y > 0 \text{ and }x + y > 0\},\\
        N &= \{(x,y) \mid x,y > 0\}.
    \end{align*}
    Let $F: M \to N$ be the function defined by $F(x,y) = \left(x + y, \dfrac{x}{y} + 1\right)$. Then $F$ is diffeomorphism since its inverse can be easily computed, we thus obtain $F^{-1}(x,y) = \left(x - \dfrac{x}{y}, \dfrac{x}{y}\right)$. Let 
    \[
        X_{(x,y)} = y^2 \frac{\pa}{\pa x}\bigg|_{(x,y)}.
    \]
    We have 
    \[
        DF(x,y) = \begin{pmatrix}
                1 & 1\\
                \dfrac{1}{y} & \dfrac{-x}{y^2}
        \end{pmatrix} \ra DF(F^{-1}(x,y)) = \begin{pmatrix}
            1 & 1\\
            \dfrac{y}{x} & \dfrac{y - y^2}{x}
        \end{pmatrix}.
    \]
    And 
    \[
        X_{F^{-1}(x,y)} = \frac{x^2}{y^2}\frac{\pa}{\pa x}\bigg|_{(x,y)}.
    \]
    It follows that 
    \[
        (F_*X)_{(x,y)} = \frac{x^2}{y^2}\frac{\pa}{\pa x}\bigg|_{(x,y)} + \frac{x}{y}\frac{\pa}{\pa y}\bigg|_{(x,y)}.
    \]
\end{example}
\section{Lie Bracket}
\begin{definition}
    Let $X,Y \in \mc{X}(M)$ be smooth vector fields. The \textbf{Lie bracket of $X$ and $Y$} is an operator $[X,Y] : C^\infty(M) \to C^\infty(M)$ defined by 
    \[
        [X,Y]f = XYf - YXf,\quad \text{for all }f \in C^\infty(M).
    \]
    The value of the bracket $[X,Y]$ at a point $p \in M$ is given by 
    \[
        [X,Y]_pf = X_p(Yf) - Y_p(Xf).
    \]
\end{definition}
\begin{lemma}
    The Lie bracket of any pair of smooth vector fields is a smooth vector field.
\end{lemma}
\begin{proof}
    Let $X,Y$ be smooth vector fields on $M$, since $[X,Y]$ is clearly linear by its definition, it is enough to show that $[X,Y]$ satisfies the Leibniz rule. For arbitrary $f,g \in C^\infty(M)$, we compute
    \begin{align*}
        [X,Y](fg) &= XY(fg) - YX(fg)\\
        &= X(gY(f) + fY(g)) - Y(gX(f) + fX(g))\\
        &= gXY(f) +Y(f)X(g) + fXY(g) + Y(g) X(f)\\
        &- gYX(f) - X(f)Y(g)- fYX(g) - X(g)Y(f)\\
        &= g[X,Y]f + f[X,Y]g.
    \end{align*}
\end{proof}
\begin{proposition}
    Let $X,Y$ be smooth vector fields on smooth manifold $M$, suppose $X = X^i \pa/\pa x^i$ and $Y = Y^j \pa/ \pa x^j$ in some coordinate chart $(x^i)$ for $M$. Then 
    \[
        [X,Y] = (XY^j - YX^j)\frac{\pa}{\pa x^j}.
    \]
\end{proposition}
\begin{proof}
    We have 
    \begin{align*}
        [X,Y]f &= X^i \frac{\pa}{\pa x^i}\left(Y^j\frac{\pa f}{\pa x^j}\right) - Y^j\frac{\pa}{\pa x^j}\left(X^i \frac{\pa f}{\pa x^i}\right)\\
        &= X^i\frac{\pa Y^j}{\pa x^i}  \frac{\pa f}{\pa x^j}+ X^iY^j\frac{\pa^2 f}{\pa x^i \pa x^j}- Y^j\frac{\pa X^i}{\pa x^j}\frac{\pa f}{\pa x^i} - X^iY^j\frac{\pa^2 f}{\pa x^j\pa x^i}\\
        &= X^i\frac{\pa Y^j}{\pa x^i}  \frac{\pa f}{\pa x^j} - Y^j\frac{\pa X^i}{\pa x^j}\frac{\pa f}{\pa x^i}\\
        &= \left(X^j\frac{\pa Y^i}{\pa x^j} - Y^j \frac{\pa X^i}{\pa x^j}\right)\frac{\pa f}{\pa x^i}\\
        &= (XY^i - YX^i)\frac{\pa}{\pa x^i}
    \end{align*}
\end{proof}
\begin{proposition}
    Let $X,Y,Z$ be smooth vector fields, then
    \begin{enumerate}
        \item Antisymmetry:
        \[
            [X,Y] = -[Y,X].
        \]
        \item Bilinearity: For $a,b \in \bb{R}$,
        \[
                [aX + bY,Z] = a[X,Z] + b[Y,Z].
        \]
        \item Jacobi identity: 
        \[
            [X,[Y,Z]] + [Y,[Z,X]] + [Z,[X,Y]] = 0.
        \]
        \item For $f,g \in C^\infty(M)$,
        \[
            [fX,gY] = fg[X,Y] + (fXg)Y - (gYf)X.
        \]
    \end{enumerate}
\end{proposition}